\begin{lemma}
In a $3$-uniform $3$-simple local extremal setup at scale $D$, let
$X_s$ be the maximal set of lowest-degree vertices in $X$ whose total degree
into $f$ is at most $D/2$, and let $X_b=X\setminus X_s$, as recorded in
\noderef{ThreeUniformThreeSimpleLocalSetup}.  Assume explicitly the standing
pair cutoff $P(f)\le1.031D^2-9D$.  If $\delta\ge0$ and
\[
P(f)\le (1+\delta)D^2-9D,
\]
then
\[
\sum_{v\in f}\sum_{x\in X_s}\deg_{\mathcal F}(v,x)
\le \frac85\delta D
\]
and
\[
w(X_s)\le \frac{2}{75}\delta D^2 .
\]
\end{lemma}
\begin{proof}
Put
\[
\varepsilon D=
\sum_{v\in f}\sum_{x\in X_s}\deg_{\mathcal F}(v,x).
\]
The definition node \noderef{ThreeUniformThreeSimpleLocalSetup} records that
$D$ is in the sufficiently large range, so $D>0$, and that the total over
$X_s$ is at most $D/2$.  Hence $\varepsilon\ge0$ and
$\varepsilon\le1/2$.

First suppose $\varepsilon=0$.  Then the total of all nonnegative table
entries in the $X_s$ columns is zero.  Each such entry is therefore zero, and
every product contributing to the weight of an $X_s$ column is zero.  Thus
$\sum_{v\in f}\sum_{x\in X_s}\deg_{\mathcal F}(v,x)=0$ and
$w(X_s)=0$, so both required inequalities follow from $\delta\ge0$.

It remains to consider $\varepsilon>0$.  The local setup gives the row and
column table facts used here: the entries are
$a_{v,x}=\deg_{\mathcal F}(v,x)$, every column sum is at most $D$, the total
sum over all of $X$ is $6(D-1)$, the total over $X_s$ is $\varepsilon D$,
and $X_b=X\setminus X_s$.  Therefore the total over $X_b$ is
$6(D-1)-\varepsilon D\le (6-\varepsilon)D$.  Applying
\noderef{TablePairFillingBound} with $k=3$, $C=D$, and
$T=(6-\varepsilon)D$ gives $q=5$ and $r=(1-\varepsilon)D$, because
$0<\varepsilon\le1/2$.  Hence
\[
w(X_b)\le \frac13\left(5D^2+(1-\varepsilon)^2D^2\right).
\]
For the $X_s$ columns, every column sum is at most the total
$\varepsilon D$, since all entries are nonnegative.  Applying the same node
to the $X_s$ columns, now with $C=T=\varepsilon D>0$, gives $q=1$ and
$r=0$, and so
\[
w(X_s)\le \frac13(\varepsilon D)^2.
\]
Since \noderef{ThreeUniformThreeSimpleLocalSetup} states
$w(X)=w(X_s)+w(X_b)$, we obtain the upper bound
\[
w(X)=w(X_b)+w(X_s)
\le \left(2-\frac23\varepsilon+\frac23\varepsilon^2\right)D^2 .
\]

The same setup records the exact pair identity
\[
P(f)=3D^2-6D+3-w(X)+Y
\]
and $Y\ge0$.  Combining this identity with the arbitrary hypothesis
$P(f)\le(1+\delta)D^2-9D$ gives
\[
w(X)=3D^2-6D+3+Y-P(f)
\ge (2-\delta)D^2+3D+3+Y
\ge (2-\delta)D^2.
\]
Comparing this lower bound with the upper bound just proved, and dividing by
the positive number $D^2$, yields
\[
\frac23\varepsilon-\frac23\varepsilon^2\le\delta .
\]

We next get a small absolute bound on $\varepsilon$ from the fixed pair
cutoff, not from the arbitrary $\delta$.  By the explicit cutoff hypothesis,
$P(f)\le1.031D^2-9D$.  Repeating the preceding rearrangement with this
bound in place of the arbitrary hypothesis gives
\[
w(X)\ge (3-1.031)D^2+3D+3+Y\ge1.969D^2.
\]
Comparing again with
$w(X)\le(2-\frac23\varepsilon+\frac23\varepsilon^2)D^2$ gives
\[
\frac23\varepsilon(1-\varepsilon)\le0.031.
\]
Since $0\le\varepsilon\le1/2$, the left side is increasing on this interval.
Equivalently, the allowed values below $1/2$ satisfy
\[
\varepsilon\le \frac12-\frac12\sqrt{1-6\cdot0.031}<0.049 .
\]
Returning to the arbitrary-$\delta$ inequality and using
$\varepsilon\le0.049$, we have
\[
\frac{2\varepsilon}{3}(1-0.049)
\le \frac23\varepsilon(1-\varepsilon)
\le \delta .
\]
Because $(2/3)(1-0.049)>5/8$, this implies
$\varepsilon\le 8\delta/5$.  Multiplying by $D$ proves
\[
\sum_{v\in f}\sum_{x\in X_s}\deg_{\mathcal F}(v,x)
=\varepsilon D\le \frac85\delta D.
\]
Finally, using the $X_s$ filling bound and the two bounds
$\varepsilon\le8\delta/5$ and $\varepsilon\le0.049$,
\[
w(X_s)\le \frac13\varepsilon^2D^2
\le \frac13\left(\frac{8\delta}{5}\right)(0.049)D^2
\le \frac{2}{75}\delta D^2 .
\]
This proves both estimates.
\end{proof}
